%=======================================================================
%
% ------------------------------------------------------------------------
% ------------------------------------------------------------------------
% abnTeX2: Modelo de Trabalho Academico  em conformidade com 
% ABNT NBR 14724:2011: Informacao e documentacao - Trabalhos academicos
% ------------------------------------------------------------------------
%Customização Dept. Engenharia Civil UFF
% ------------------------------------------------------------------------
% Autor:  Malú Grave (adaptado de Jonas Alessi (alessi.jonas@gmail.com) - UNOESC, 2013)
% Versão: 2 de Abril 2023.
% Edição: TexStudio
% Codificação: UTF-8
% LaTeX:  abnTeX2
%
%=======================================================================

\documentclass[
	% -- opções da classe memoir --
	12pt,				% tamanho da fonte	
	oneside,          % não imprimir em verso e anverso, oposto do twoside 
	a4paper,			% tamanho do papel. 
	% -- opções da classe abntex2 --
	chapter=TITLE,		% títulos de capítulos convertidos em letras maiúsculas
	section=TITLE,		% títulos de seções convertidos em letras maiúsculas
	subsection=TITLE,	% títulos de subseções convertidos em letras maiúsculas
	%subsubsection=TITLE,% títulos de subsubseções convertidos em letras maiúsculas
	% -- opções do pacote babel --
	english,			% idioma adicional para hifenização
	brazil			% o último idioma é o principal do documento
	,sumario=tradicional
	]{abntex2-uff}


\include{config_inicial}

%----Include da capa é fora do documento 
\include{pretextuais/capa}
%---

\begin{document}
% Retira espaço extra obsoleto entre as frases.
\frenchspacing 

% ----------------------------------------------------------
% ELEMENTOS PRÉ-TEXTUAIS
% ----------------------------------------------------------

%--- CAPA ----
\imprimircapa

% --- FOLHA DE ROSTO
\imprimirfolhaderosto

% Inserir ficha catalográfica
\includepdf[pages=-]{pretextuais/ficha.pdf}
%Ficha Catalográfica elaborada pelo Aluno, acesse o site: https://bibliotecas.uff.br/bcg/fichacatalografica/

%Gere o pdf da ficha e substitua o arquivo ficha.pdf com a sua ficha.

% Inserir FOLHA DE APROVAÇÃO
\include{pretextuais/folhadeaprovacao}

% DEDICATÓRIA
\include{pretextuais/dedicatoria}

% AGRADECIMENTO
\include{pretextuais/agradecimentos}

%EPÍGRAFE
%\include{pretextuais/epigrafe}

% RESUMO
\include{pretextuais/resumos}

% lista de ilustrações
\pdfbookmark[0]{\listfigurename}{lof}
\listoffigures*
\cleardoublepage
% ---


% LISTA DE TABELAS
\pdfbookmark[0]{\listtablename}{lot}
\listoftables*
\cleardoublepage  %-- força proxima pagina

% ---
% inserir lista de abreviaturas e siglas
% ---
\begin{siglas}
  \item[ABNT] Associação Brasileira de Normas Técnicas
  \item[Bras.] Brasileiro
  \item[CNPq] Conselho Nacional de Desenvolvimento Científico e Tecnológico
  \item[IBGE] Instituto Brasileiro de Geografia e Estatística
  \item[LCD] \textit{Liquid Crystal Dysplay}
\end{siglas}
% ---

% ---
% inserir lista de símbolos
% ---
\begin{simbolos}
  \item[$ \Gamma $] Letra grega Gama
  \item[$ \Lambda $] Lambda
  \item[$ \zeta $] Letra grega minúscula zeta
  \item[$ \in $] Pertence
\end{simbolos}
% ---

% SUMARIO
\pdfbookmark[0]{\contentsname}{toc}
\tableofcontents*
\cleardoublepage

% ------------------------------------------------------
% ELEMENTOS TEXTUAIS
% ------------------------------------------------------
\textual
% INTRODUÇÃO
\include{conteudo/introducao}

%FUNDAMENTAÇÃO TEÓRICA
\include{conteudo/fundamentacao}

%DESENVOLVIMENTO
\include{conteudo/desenvolvimento}

%RESULTADOS
\include{conteudo/resultados}

%CONCLUSÃO
\include{conteudo/conclusao}

% ----------------------------------------------------------
% ELEMENTOS PÓS-TEXTUAIS
% ----------------------------------------------------------
\postextual

% ----------------------------------------------------------
% Referências bibliográficas
% ----------------------------------------------------------

\bibliography{referencias}

%Apêndices
\include{postextuais/apendice}

%Anexos
\include{postextuais/anexo}

%---------------------------------------------------------------------
% INDICE REMISSIVO
%---------------------------------------------------------------------

\printindex
\end{document}